% Preambolo comune dell'allegato tecnico. Si compila con XeLaTeX + BibTeX.
\documentclass[10pt,a4paper]{article}
\usepackage[a4paper,margin=2.2cm]{geometry}
\usepackage{fontspec}
\setmainfont{Fira Sans}[Scale=0.95,BoldFont={Fira Sans SemiBold},ItalicFont={Fira Sans Italic}]
\setmonofont{DejaVu Sans Mono}[Scale=0.82]
\usepackage{polyglossia}\setmainlanguage{italian}
\usepackage{amsmath,amssymb,amsthm,mathtools}
\usepackage{xcolor,tikz,booktabs,tabularx,array,enumitem,microtype,ragged2e,caption,float}
\usetikzlibrary{positioning,arrows.meta,calc,angles,quotes,patterns,decorations.pathreplacing}
\usepackage[numbers,sort&compress]{natbib}
\usepackage[hidelinks]{hyperref}
\definecolor{accent}{HTML}{0B5C8A}
\definecolor{muted}{HTML}{55606C}
\definecolor{boxbg}{HTML}{EEF3F7}
\setlength{\parindent}{0pt}\setlength{\parskip}{4pt}
\setlist{leftmargin=1.4em,itemsep=1.5pt,topsep=2pt}
\captionsetup{font=small,labelfont=bf}
\newcolumntype{L}{>{\RaggedRight\arraybackslash}X}
% Ambienti numerati per capitolo (sezione)
\theoremstyle{plain}
\newtheorem{proposizione}{Proposizione}[section]
\newtheorem{lemma}[proposizione]{Lemma}
\newtheorem{corollario}[proposizione]{Corollario}
\theoremstyle{definition}
\newtheorem{definizione}[proposizione]{Definizione}
\newtheorem{ipotesi}[proposizione]{Ipotesi}
\newtheorem{controesempio}[proposizione]{Controesempio}
\theoremstyle{remark}
\newtheorem{osservazione}[proposizione]{Osservazione}
% Livello di prova di un risultato: \prova{D} dimostrato, {C} calcolo, {S} simulazione, {P} dati pubblici, {N} non verificato sperimentalmente
\newcommand{\prova}[1]{\textsf{\footnotesize\color{accent}[#1]}}
% Riquadro per condizioni di validità e limiti
\usepackage{framed}\colorlet{shadecolor}{boxbg}
\newenvironment{riquadro}[1]{\begin{shaded}\textbf{#1}\par\smallskip}{\end{shaded}}
% Notazione comune
\newcommand{\Reali}{\mathbb{R}}
\newcommand{\vmax}{v_{\max}}
\newcommand{\tocc}{\tau_{\mathrm{occ}}}
\newcommand{\Oss}{\mathcal{O}}      % insieme osservato libero
\newcommand{\Omb}{\Omega}           % ombra
\newcommand{\Pres}{\Pi}             % possibile presenza
\newcommand{\Reg}{\mathcal{R}}      % regione di previsione
\newcommand{\unita}[1]{\,\mathrm{#1}}
\DeclareMathOperator{\dilata}{dil}

\begin{document}
\paragraph{Frammento 1.}
\begin{proposizione}[Contenimento]\label{prop:occ-contenimento}
Sia $\Oss(t)$ l'insieme dei punti del suolo il cui raggio, nel fotogramma $t$, è letto da un pixel valido alla profondità attesa del suolo (lettura per punto \emph{esatta}: margine di lettura nullo nella \eqref{eq:map-margine}). Sotto le ipotesi (i), (iv) e (v), a posa esatta ($\varepsilon_p=0$), ogni punto interno all'impronta al suolo di un oggetto al fotogramma $t$ (punto \emph{occupato}) appartiene a $\Pres(t)$. In particolare la posizione di un soggetto entrato in un'ombra appartiene a $\Pres(t)$ per ogni $t$ fino alla riosservazione, senza alcuna ipotesi sul numero dei soggetti né sull'esistenza di una traccia. Con posa nota entro $\varepsilon_p$ (ipotesi (iii), intesa come spostamento nel piano del suolo, come nel Lemma~\ref{lem:map-posa}) l'enunciato vale nel riferimento stimato; nel riferimento vero vale se nella \eqref{eq:occ-presenza} a $\Oss(t)$ si sostituisce l'insieme dei punti $x$ il cui disco di raggio $\varepsilon_p$ è contenuto nel libero stimato $\hat{\Oss}(t)$. \prova{D}
\end{proposizione}

\paragraph{Frammento 2.}
Per (i) e (v) un raggio letto alla profondità attesa del suolo è privo di superfici fino al suolo; il raggio diretto a un punto interno all'impronta incontra invece l'oggetto prima del suolo, quindi un punto occupato non è letto come suolo: $x\notin\Oss(t)$.

\paragraph{Frammento 3.}
Con il margine di lettura $m>0$ della \eqref{eq:map-margine} il passo «un punto occupato non è letto come suolo» cade su una striscia dell'impronta. Il raggio diretto a un punto che sta dietro la faccia vista dell'oggetto, a distanza $\delta$ in pianta, incontra la faccia alla profondità $z_0\,(1-\delta/\rho)$: se $\delta\le\rho\,m(z_0)/z_0$ la lettura dista da $z_0$ meno del margine e il punto è letto come suolo, qualunque sia l'altezza dell'oggetto. Sull'asse del campo, con $h=1{,}2\unita{m}$, $\theta=15^\circ$ e $k=0{,}004$, la striscia è larga $11\unita{cm}$ a $2{,}5\unita{m}$ e $14\unita{cm}$ a $3\unita{m}$; con $k=0{,}0100$, $23$ e $32\unita{cm}$ \prova{C}. La Proposizione è quindi dimostrata per la lettura per punto esatta \prova{D}; con il margine di lettura la parte dell'impronta più vicina alla camera può uscire da $\Pres$, e il contenimento dell'intera impronta non è dimostrato. Nemmeno per la regola della cella (Definizione~\ref{def:map-libero}) il passo è dimostrato: una cella può essere «libero» con meno di metà dell'impronta coperta da un oggetto. Con il margine di lettura e con la mappa a celle il contenimento poggia sulla verifica in simulazione del \S\ref{sec:map-banco} \prova{S}.

\paragraph{Frammento 4.}
Sotto le ipotesi (ii) e (iv), con la (ii) intesa per il centro (un soggetto con il centro fuori da $\Omb$, in campo e in portata, è riportato come osservato entro $t_a$), se $w\ge0$.

\paragraph{Frammento 5.}
Per contenere l'intero corpo di raggio $r_c$ si aggiunge $r_c$ al secondo membro; se la (ii) vale solo quando tutto il corpo è fuori dall'ombra, se ne aggiunge un altro. \prova{D}

\paragraph{Frammento 6.}
Con $P_D<1$ nessun tempo finito è un limite certo: per una traccia nuova la conferma entro due cicli ha probabilità $0{,}81$ con $P_D=0{,}9$ (\S\ref{sec:sti-porta}).

\paragraph{Frammento 7.}
\begin{tabular}{lp{10cm}}
$t_a$ & tempo entro cui un soggetto fuori dall'ombra, in campo e in portata, è riportato come osservato\\
\end{tabular}

\paragraph{Frammento 8.}
Lo zero della prima riga vale per costruzione dove l'ingresso è riconosciuto, l'ostacolo è fermo e l'ombra stimata coincide con la vera: allora $W(t)=\dilata(\hat\Omb,w)\supseteq\Omb$ (Definizione~\ref{def:occ-avviso}) e il centro di un soggetto nascosto sta in $\Omb$; altrove è un conteggio. Non è un'applicazione della Proposizione~\ref{prop:occ-fascia}, che riguarda il soggetto uscito dall'ombra e non ancora riportato come osservato, e la cui condizione \eqref{eq:occ-fascia} copre $1{,}13\unita{m/s}$, meno delle velocità d'ingresso di queste prove (fino a $1{,}8\unita{m/s}$).

\paragraph{Frammento 9.}
Un'etichetta non produce un «libero» e, per ogni ingresso in un'ombra riconosciuto, non riduce l'area d'avviso (Proposizione~\ref{prop:occ-etichetta}); può cambiare l'esito di un riaggancio, e quindi gli ingressi successivi (\S\ref{sec:occ-classe}).

\paragraph{Frammento 10.}
Un'etichetta non produce un «libero» e non toglie né restringe un avviso acceso (invarianti 5 e 15); può cambiare l'esito di un riaggancio, e quindi gli ingressi successivi (\S\ref{sec:occ-classe}).

\paragraph{Frammento 11.}
non entra nel costo d'associazione né nella porta di validazione; entra invece, attraverso la forma della regione di previsione, nella porta del riaggancio (\S\ref{sec:occ-classe});

\paragraph{Frammento 12.}
$\Reg_\alpha(t)=C_\alpha(t)\cap W(t)$ per $\alpha<1$; se questo insieme è vuoto, e in ogni altro caso di «previsto ampio», $\Reg_\alpha(t)=\bigl(\hat\Omb(t)\setminus\Oss(t)\bigr)\cap W(t)$, cioè tutto il non visto libero dell'ombra (\emph{regione di ripiego}). Si pone $\Reg_1(t)=W(t)$.

\paragraph{Frammento 13.}
(d) $\hat\Omb\cap\Oss(t)=\emptyset$ per definizione d'ombra e, con l'ostacolo fermo e l'ombra stimata costante, $\hat\Omb\subseteq W(t_0)=W(t)$; quindi la regione di ripiego è $\hat\Omb\neq\emptyset$.

\paragraph{Frammento 14.}
parte (b) dimostrata nel modello a posa esatta; con errore d'inclinazione entro $0{,}5^\circ$ la tenuta è verificata in simulazione, non dimostrata;

\paragraph{Frammento 15.}
l'ostacolo è fermo, il soggetto rimanda segnale, resta nel campo e in portata ed è riportato come osservato entro $t_a$ da quando il suo centro è fuori dall'ombra
\end{document}
